\newpage
\section*{Wprowadzenie}

Uczenie maszynowe niejednokrotnie w różnych gałęziach przemysłu jest wykorzystywane do analizy skomplikowanych i nieliniowo powiązanych ze sobą danych. \cite{review_predictive_maintenance_2023} Jednym z powszechnych zastosowań jest analiza przeżywalności komponentów i całych silników lotniczych \cite{saxena2008damage}, dzięki której producenci silników (np. General Electric) mogą oszacować liczbę napraw w następującym roku dla każdego silnika, która determinuje przychody z tytułu utrzymywania zdatności do lotu floty klienta (linii lotniczej). Te szacowane przychody następnie są raportowane akcjonariuszom co pośrednio wpływa na wycenę akcji. 

Skutkiem zniszczeń maszyn w przemyśle może być strata finansowa lub uszczerbek na zdrowiu bądź strata życia ludzi. Wczesne wykrywanie uszkodzeń poza minimalizacją kosztów eksploatacji i wymiany urządzeń jest nieodzowne dla bezpiecznego i ciągłego funkcjonowania. Z tego powodu implementacja uczenia maszynowego, a w szczególności autoenkoderów do szacowania RUL (ang. Remaining Useful Life, tj. dosł. Pozostałego Użytecznego Życia) jest istotna i skutecznie wykorzystywana. \cite{autoencoder_fault_detection_2022}

Wspomniana skuteczność sztucznych sieci neuronowych zarówno klasycznych jak i kwantowych wynika z ich zdolności do modelowania nieliniowych zależności pomiędzy danymi. Nieliniowość i mnogość danych (silniki mogą implementować kilka tysięcy czujników pomiarowych generujących dziesiątki GB danych po każdym locie) prognostycznych silników lotniczych jest jednym z głównych czynników, które utrudniają szacowanie RUL. Jest to skutek złożonych i wzajemnie sprzężonych interakcji aerodynamicznych i termodynamicznych zachodzących w~silniku. Tradycyjne metody statystyczne nie są w stanie modelować tych zjawisk z wystarczającą dokładnością.

Autoenkodery wyróżniają się na tle innych narzędzi takich jak Analiza Głównych Składowych (ang. Principal Component Analysis) wyższą jakością rekonstrukcji oraz kompresji danych (zmniejszania wymiarowości). Kompresja jest kluczowa ze względu na nieustannie rosnący popyt na składowanie danych (ang. Data storage) wraz z podażą, która nie rośnie tak szybko jak popyt. Ponadto zmniejszanie wymiarowości sprzyja wizualizacji oraz komunikacji oraz przyspiesza analizę. \cite{hinton2006reducing}

Obliczenia kwantowe są w stanie modelować zjawiska, które przekraczają modelowanie oparte na fizyce klasycznej. Innymi słowy systemy uczenia maszynowego oparte na rozwiązaniach kwantowych przekraczają możliwościami rozpoznawania wzorców klasyczne systemy. \cite{romero2017quantum} Wraz z rozwojem technologii kwantowych oraz ze słabnącym prawem Moore'a (ang. Moore's Law) systemy kwantowe stają się coraz bardziej atrakcyjną alternatywą dla klasycznych systemów komputerowych, lub metodą by usprawnić ich działanie.

Celem niniejszej pracy jest opracowanie modeli klasycznych sztucznych sieci neuronowych oraz kwantowych sieci neuronowych oraz ich analiza. Wykorzystane zostaną w celu kompresji oraz dekompresji danych diagnostycznych silników turbowentylatorowych.

Główną motywacją badawczą jest optymalizacja procesów telemetrii lotniczej. Implementacja wyuczonego enkodera na pokładzie statku powietrznego umożliwia minimalizację strumienia przesyłanych danych oraz zmniejszenie zapotrzebowania na fizyczną przestrzeń dyskową systemów pokładowych. W kontekście lotnictwa cywilnego redukcja masy komponentów elektronicznych ma kluczowe znaczenie dla efektywności paliwowej i tym samym kosztów operacyjnych.

W pracy podjęto próbę oceny, jak charakterystyka danych wejściowych – w tym ich nieliniowość, nieciągłość oraz rozkład statystyczny – wpływa na jakość rekonstrukcji sygnału przez modele klasyczne i kwantowe. Badania przeprowadzono z wykorzystaniem benchmarkowego zbioru danych C-MAPSS (ang. Commercial Modular Aero-Propulsion System Simulation), opracowanego przez centrum badawcze NASA \cite{saxena2008damage}.